\section{Method}
\subsection{Threat Model and Fast Update}
Let $f_\theta$ be a classifier and $\ell$ cross-entropy loss. Adversarial
training under an $\ell_\infty$ radius $\epsilon$ approximates
\begin{equation}
\min_\theta\;\mathbb{E}_{(x,y)}
\left[\max_{\|\delta\|_\infty\leq\epsilon}
\ell(f_\theta(x+\delta),y)\right].
\end{equation}
FGSM begins at $x^0$, either $x$ or a random point in the perturbation set,
and constructs
\begin{equation}
g^0=\nabla_{x^0}\ell(f_\theta(x^0),y),\quad
x^{\rm adv}=\Pi\left(x^0+\alpha\,\mathrm{sign}(g^0)\right),
\end{equation}
where $\Pi$ projects onto the valid image and threat set.

\subsection{Attack-Free Gradient-Consistency Score}
The endpoint input gradient
$g^1=\nabla_{x^{\rm adv}}\ell(f_\theta(x^{\rm adv}),y)$ is obtained from the
same parameter-backward pass, without an additional forward or backward pass.
We retain it and calculate
\begin{equation}
c(x)=\frac{\langle g^0,g^1\rangle}
{\|g^0\|_2\|g^1\|_2}.
\end{equation}
FGSM relies on a local first-order approximation. A sharp decrease in the
cosine between $g^0$ and $g^1$ means that the input gradient rotates along the
FGSM direction, so the local linear approximation and the representativeness
of the single-step inner maximizer are degrading. This interpretation agrees
with prior links between local nonlinearity and CO
\cite{andriushchenko2020understanding}. We do not claim that the cosine drop is
a necessary or sufficient condition for CO; it is an empirically validated
proxy for the local-linearity failure associated with CO.

Batch averages are accumulated over non-overlapping windows of 20 optimizer
updates to give $c_t$. The monitor adds reductions and trace state, but no
additional forward or backward pass. In the implementation, $g^0$ constructs
the FGSM example and $g^1$ is retained from the same parameter-backward pass;
the monitor calls no additional autograd-gradient operation.
To remove the slow training trend, it uses
\begin{equation}
b_t=\mathrm{median}(c_{t-5},\ldots,c_{t-1}),\qquad s_t=b_t-c_t,
\end{equation}
after at least three history windows. A warning occurs when $s_t\geq\tau$.
For CIFAR-10 calibration, we pool every finite development score, label a
window positive only when a registered CO event is the next 20-update window,
scan the unique observed scores, and choose the threshold maximizing Youden's
$J$ (true-positive rate minus false-positive rate). Seeds 17, 23, and 42 form
the development set; this yields $\tau=0.1663298875$, which is then frozen.
CIFAR-100 uses a separate event-level lexicographic rule, detailed
in Section IV. Offline sensitivity analysis in the Appendix does not alter
either fitted value or the main protocol.

\subsection{Immediate-Switch Conditional Escalation}
At the first warning, the primary response retains the completed FGSM update,
switches the next mini-batch to PGD-2, and uses PGD-2 for all subsequent
updates. Thus no additional FGSM update intervenes between warning and action,
and no training state must be restored.

We call this no-replay main response \emph{immediate-switch}. Four matched
alternatives are evaluated. The rollback ablation, \emph{immediate-replay},
keeps an in-memory epoch-start snapshot of model, optimizer, mixed-precision
scaler, Python, NumPy, PyTorch CPU/CUDA, and loader-generator random states.
At the warning it restores that state and replays the epoch using PGD-2; the
snapshot is not serialized to disk. The remaining controls retain the current
FGSM epoch and start PGD-2 at the next epoch, switch at a fixed early or late
epoch, or use PGD-2 from epoch 1. These controls separate the value of prompt
action from any effect of restoring the epoch.

\subsection{Event and Warning/Evaluation Protocol}
On fixed validation data, a CO event requires PGD-10 accuracy to fall at least
20 percentage points from its previous peak and reach at most 10\%. We report
whether a warning precedes the first event, whether it occurs within 60
updates, median lead, and alarm episodes in non-event runs. Validation PGD is
used for retrospective event labels and for selecting \emph{best\_robust.pt};
it never triggers the proposed monitor. Thus ``attack-free'' describes the
online monitor, not an attack-free checkpoint-selection protocol. Recall
intervals are Wilson 95\% intervals. Accuracy and time comparisons use
identical seeds and paired bootstrap 95\% intervals with 100,000 resamples.
